\documentclass[varwidth=16cm, border=2mm]{standalone}
\usepackage{amsmath}\usepackage{amssymb}\usepackage{ctex}
\usepackage{tasks}\usepackage{xcolor}\usepackage{graphicx}
\settasks{label=\Alph*., label-width=1.5em, item-indent=2em}
\begin{document}
\textbf{[0.6]} （2024·全国·高三对口高考）已知椭圆 $\frac{x^2}{9} + y^2 = 1$，过左焦点 $F$ 作倾斜角为 $\frac{\pi}{6}$ 的直线交椭圆于 $A$、$B$ 两点，则弦 $AB$ 的长为\underline{\hspace{1.5cm}}.

\vspace{0.2cm}\par\noindent\textcolor{red}{\textbf{【答案】} 2}
\par\noindent\textcolor{blue}{\textbf{【解析】} 已知椭圆方程为 $\frac{x^2}{9} + y^2 = 1$。
由此可知 $a^2 = 9$，所以 $a = 3$；$b^2 = 1$，所以 $b = 1$。
椭圆的半焦距 $c$ 满足 $c^2 = a^2 - b^2 = 9 - 1 = 8$，所以 $c = \sqrt{8} = 2\sqrt{2}$。
左焦点 $F$ 的坐标为 $(-c, 0) = (-2\sqrt{2}, 0)$。
直线 $AB$ 的倾斜角为 $\frac{\pi}{6}$，所以直线的斜率 $k = \tan(\frac{\pi}{6}) = \frac{\sqrt{3}}{3}$。

因此，过左焦点 $F(-2\sqrt{2}, 0)$ 且倾斜角为 $\frac{\pi}{6}$ 的直线 $l$ 的方程为：
$y - 0 = \frac{\sqrt{3}}{3}(x - (-2\sqrt{2}))$
$y = \frac{\sqrt{3}}{3}(x + 2\sqrt{2})$

将直线方程代入椭圆方程 $\frac{x^2}{9} + y^2 = 1$：
$\frac{x^2}{9} + (\frac{\sqrt{3}}{3}(x + 2\sqrt{2}))^2 = 1$
$\frac{x^2}{9} + \frac{3}{9}(x + 2\sqrt{2})^2 = 1$
$\frac{x^2}{9} + \frac{1}{3}(x^2 + 4\sqrt{2}x + 8) = 1$
等式两边同乘以 9，得：
$x^2 + 3(x^2 + 4\sqrt{2}x + 8) = 9$
$x^2 + 3x^2 + 12\sqrt{2}x + 24 = 9$
$4x^2 + 12\sqrt{2}x + 15 = 0$

设直线与椭圆的两个交点为 $A(x_1, y_1)$ 和 $B(x_2, y_2)$。
根据韦达定理，我们有：
$x_1 + x_2 = -\frac{12\sqrt{2}}{4} = -3\sqrt{2}$
$x_1 x_2 = \frac{15}{4}$

弦 $AB$ 的长度可以使用弦长公式 $L = \sqrt{1+k^2}|x_1-x_2|$ 进行计算。
首先计算 $|x_1-x_2|$：
$|x_1-x_2|^2 = (x_1+x_2)^2 - 4x_1x_2$
$= (-3\sqrt{2})^2 - 4(\frac{15}{4})$
$= (9 \times 2) - 15$
$= 18 - 15 = 3$
所以 $|x_1-x_2| = \sqrt{3}$。

$k = \frac{\sqrt{3}}{3}$，所以 $k^2 = (\frac{\sqrt{3}}{3})^2 = \frac{3}{9} = \frac{1}{3}$。

弦 $AB$ 的长 $L = \sqrt{1+k^2}|x_1-x_2| = \sqrt{1 + \frac{1}{3}} \times \sqrt{3}$
$= \sqrt{\frac{4}{3}} \times \sqrt{3}$
$= \frac{2}{\sqrt{3}} \times \sqrt{3}$
$= 2$。

也可以使用焦半径公式。对于椭圆 $\frac{x^2}{a^2} + \frac{y^2}{b^2} = 1$，其离心率 $e = \frac{c}{a} = \frac{2\sqrt{2}}{3}$。
过左焦点 $F$ 的弦 $AB$ 的长度为 $AB = AF + BF$ (此为焦点弦的性质)。
对于椭圆上的点 $P(x,y)$，它到左焦点 $F$ 的距离为 $PF = a + ex$。
所以 $AF = a + ex_1$，$BF = a + ex_2$。
弦 $AB$ 的长为 $AF + BF = (a + ex_1) + (a + ex_2) = 2a + e(x_1+x_2)$。
将 $a=3$, $e=\frac{2\sqrt{2}}{3}$, $x_1+x_2=-3\sqrt{2}$ 代入：
$AB = 2(3) + \frac{2\sqrt{2}}{3}(-3\sqrt{2})$
$= 6 + \frac{-12}{3}$
$= 6 - 4 = 2$。

最终答案为 2。}


\end{document}
